\section{Zadanie 2}
Jeśli \({A_n}_{n\in \mathbb{N}}\) jest występującą rodziną zdarzeń, czyli
\(A_0 \subseteq A_1 \subseteq A_2 \subseteq \ldots\), to:
\[
	P\left(\bigcup_{n=0}^{\infty} A_n\right) = \lim_{n \to \infty} P(A_n)
\]

Jeśli \({A_n}_{n\in \mathbb{N}}\) jest zstępującą rodziną zdarzeń, czyli
\(A_0 \supseteq A_1 \supseteq A_2 \supseteq \ldots\), to:
\[
	P\left(\bigcap_{n=0}^{\infty} A_n\right) = \lim_{n \to \infty} P(A_n)
\]


\subsection*{Rozwiązanie}
Tak to prawda XD;

\subsubsection*{Dowód}
Poniżej podamy dowody dla obu stwierdzeń. W obu przypadkach skorzystamy z prostych własności miary prawdopodobieństwa oraz definicji granicy ciągu liczb rzeczywistych (epsilon -- delta, tu: epsilon i indeks \(N\)).

\paragraph{1) Rodzina rosnąca: \(A_0\subseteq A_1\subseteq\ldots\)}

Niech \(A=\bigcup_{n=0}^{\infty}A_n\). Wykazujemy,
\[P(A)=\lim_{n\to\infty}P(A_n).\]

Ponieważ \(A_n\subseteq A_{n+1}\) dla każdego \(n\), ciąg \(P(A_n)\) jest niemalejący i ograniczony z góry przez \(1\), a zatem ma granicę (można to przyjmować jako \(L:=\lim_{n\to\infty}P(A_n)\) lub jako \(\sup_n P(A_n)\)). Pokażemy, że \(P(A)=L\).

Zauważmy, że zbior \(A\) jest sumą rozłącznych przyrostów:
\[A=\bigcup_{n=0}^{\infty}A_n= A_0\uplus (A_1\setminus A_0)\uplus (A_2\setminus A_1)\uplus\cdots,\]
gdzie symbolem \(\uplus\) oznaczamy rozłączność w sensie rozłożenia na zbiory parami rozłączne.

Dla dowolnego \(m\) mamy
\[A_m = A_0\uplus (A_1\setminus A_0)\uplus\cdots\uplus (A_m\setminus A_{m-1}).\]
Korzystając z addytywności prawdopodobieństwa na zbiorach parami rozłącznych dostajemy
\[P(A_m)=P(A_0)+\sum_{k=1}^m P(A_k\setminus A_{k-1}).\]
Jeśliśmy dla \(m\to\infty\), to z monotoniczności i jednoczesnej nieujemności wyrazów szeregu mamy
\[\lim_{m\to\infty}P(A_m)=P(A_0)+\sum_{k=1}^{\infty} P(A_k\setminus A_{k-1}).\]
Równocześnie
\[P(A)=P\Big(\biguplus_{k=0}^{\infty} (A_k\setminus A_{k-1})\Big)=P(A_0)+\sum_{k=1}^{\infty} P(A_k\setminus A_{k-1}),\]
stąd \(P(A)=\lim_{n\to\infty}P(A_n)\).

Można to również sformalizować bez rozkładu na przyrosty, używając definicji granicy: dla danego \(\varepsilon>0\) niech \(L=\lim_{n}P(A_n)=\sup_n P(A_n)\). Istnieje indeks \(N\) taki, że dla \(n\ge N\) mamy \(P(A_n)>L-\varepsilon\). Ponieważ \(A_n\uparrow A\), to \(A\setminus A_n\downarrow\emptyset\) i z podzielności miary \(P(A)-P(A_n)=P(A\setminus A_n)\to0\), skąd \(P(A)=L\).

\paragraph{2) Rodzina malejąca: \(A_0\supseteq A_1\supseteq A_2\supseteq\ldots\)}

Niech \(A=\bigcap_{n=0}^{\infty}A_n\). Zakładamy dodatkowo, że istnieje \(k\) takie, że \(P(A_k)<\infty\) (w kontekście probabilistycznym zwykle \(P\) jest miarą skończoną i ta hipoteza jest automatycznie spełniona). Chcemy pokazać, że
\[P(A)=\lim_{n\to\infty}P(A_n).\]

Rozważmy zbiory \(B_n:=A_0\setminus A_n\). Ponieważ \(A_n\) jest malejący, to \(B_n\) jest rosnący i
\[\bigcup_{n=0}^{\infty}B_n = A_0\setminus\bigcap_{n=0}^{\infty}A_n = A_0\setminus A.\]
Z pierwszej części dowodu (ciągi rosnące) mamy
\[P\Big(\bigcup_{n=0}^{\infty}B_n\Big)=\lim_{n\to\infty}P(B_n).
\]
Ale \(P(B_n)=P(A_0)-P(A_n)\) dla każdego \(n\) (ponieważ \(B_n\) i \(A_n\) rozdzielają \(A_0\)), stąd
\[P(A_0)-P(A)=\lim_{n\to\infty}\bigl(P(A_0)-P(A_n)\bigr)=P(A_0)-\lim_{n\to\infty}P(A_n).
\]
Po uproszczeniu otrzymujemy \(P(A)=\lim_{n\to\infty}P(A_n)\), co kończy dowód.

\paragraph{Uwaga o hipotezie skończoności:} Jeśli \(P(A_0)=\infty\) (nie jest to sytuacja typowa dla prawdopodobieństwa, ale może wystąpić dla ogólnych miar), to równanie \(P(B_n)=P(A_0)-P(A_n)\) nie ma sensu; w teorii miar stosuje się wtedy warunek, że istnieje przynajmniej jeden indeks \(k\) z \(P(A_k)<\infty\), co umożliwia analogiczny argument na przesunietym ciągu.



